\documentclass[10pt,twocolumn]{article}

\usepackage[margin=0.72in,columnsep=0.24in]{geometry}
\usepackage{amsmath,amssymb}
\usepackage{booktabs}
\usepackage{enumitem}
\usepackage{graphicx}
\usepackage{microtype}
\usepackage{tabularx}
\usepackage{xcolor}
\usepackage{xspace}
\usepackage{url}
\usepackage[hidelinks]{hyperref}

\newcommand{\system}{\textsc{Tempo}\xspace}
\newcommand{\slogood}{\textsc{SLO-Good}\xspace}
\newcommand{\todo}[1]{\textcolor{red}{[TODO: #1]}}
\newcommand{\tightparagraph}[1]{\vspace{0.25em}\noindent\textbf{#1.}}
\setlist{nosep,leftmargin=*}
\setlength{\parskip}{0.15em}

\title{\system: Cross-Layer Global Orchestration for\\
Disaggregated LLM Inference on Shared HPC Fabrics}
\author{Anonymous Author(s)}
\date{}

\begin{document}
\maketitle

\begin{abstract}
Prefill/decode disaggregation removes one form of GPU interference but creates
a coupled control problem: a local prefill consumes decode capacity, whereas a
remote prefill consumes KV-transfer, receiver, and network capacity.  The
better path therefore changes with workload, cache residency, decoder service,
fabric contention, and tenant objectives.  Existing inference schedulers and
network transports each expose valuable parts of this state, but optimizing one
layer can move the bottleneck to another.

We present \system, a global orchestrator that joins vLLM service state,
LMCache KV semantics, Cassini/Slingshot endpoint telemetry, topology, and
business SLO state in an identity-bound control loop.  It atomically admits a
request, selects a prefill/decode edge and local or remote execution, reserves
decoder and transfer resources, and records fail-closed lifecycle evidence.
Near-tied remote replicas are balanced by controller-owned source/edge service,
and a node--pair--shard hierarchy bounds global fan-in.

On four Perlmutter nodes (16 A100 GPUs), using Qwen2.5-7B-Instruct, actual
vLLM, LMCacheConnectorV1, NIXL/UCX over CXI, and a fresh held-out workload,
\system attains all 210 offered foreground SLOs across normal, decoder-hot, and
remote-favorable regimes.  Relative to the strongest fixed policy, it reduces
stressed p99 latency by 66.1--90.6\%; relative to a predictor, by 61.7--93.4\%.
In a same-allocation post-hoc extension, \system reduces p99 by 75.8--93.8\%
over a NetKV Algorithm-1 reproduction while improving stressed SLO attainment
from 73/120 and 0/30 to 120/120 and 30/30.  A 1,024-pair control-plane replay
reduces global payload by 87.5\%.  We explicitly separate the independently
validated four-node result from the post-hoc SOTA comparison and from
CPU-only scale evidence.
\end{abstract}

\section{Introduction}

Large-language-model serving separates prompt prefill from autoregressive
decode because their compute and memory behavior differ.  Systems such as
DistServe and Mooncake isolate the phases and move KV state between pools,
improving resource planning and cache reuse~\cite{distserve,mooncake}.  vLLM
provides the high-throughput execution substrate~\cite{vllm}, and LMCache
exposes KV state as a cross-engine storage and communication layer~\cite{lmcache}.

Disaggregation, however, does not eliminate contention.  It changes where
contention appears.  Sending a request through a prefill node can queue behind
long prompts and then pay an inter-node KV transfer.  Deflecting prefill to a
decode node avoids that transfer but competes with active decodes.  Choosing a
different decode receiver can relieve its GPU queue while creating incast or
moving load to a congested fabric edge.  In a shared HPC environment, inference
also coexists with collective communication and jobs whose utility differs by
tenant, deadline, and fairness obligation.

This creates a systems question broader than route prediction:

\begin{quote}
Can one controller combine application semantics, GPU service, KV movement,
fabric state, topology, and business demand, then jointly control admission,
placement, and execution without simply moving the bottleneck?
\end{quote}

Network-aware selection, as formalized by NetKV~\cite{netkv}, makes transfer
cost visible.  Kairos~\cite{kairos} instead exploits decode-side compute slack
with a TBT-safe chunk schedule.  Multipath transports such as MRC improve path
utilization and recovery~\cite{mrc}.  These are complementary mechanisms:
none alone owns the complete request lifecycle from tenant admission through
prefill/decode placement, KV transfer, receiver service, and completion.

\system treats this cross-layer lifecycle as the unit of control.  Its design
is guided by three observations from our development history.  First, local and
remote winners reverse as decoder and transfer pressure change.  Second,
completed-only latency is misleading when a policy rejects or loses requests.
Third, telemetry is useful only when its identity, freshness, support status,
and actuation are bound to the request that consumed it.

This paper makes four contributions:

\begin{enumerate}
  \item We formulate disaggregated inference on shared HPC fabrics as a
  cross-layer global admission and edge-selection problem, including tenant
  utility and failure state rather than latency prediction alone.
  \item We implement an identity-bound state and actuation plane spanning
  vLLM, LMCache/NIXL, Cassini/Slingshot, topology, and business policy.  The
  same transaction selects local/remote execution and physical $P_i$--$D_j$
  placement, reserves resources, and closes the stream lifecycle.
  \item We construct realistic matched workloads that separately expose
  decoder-hot and remote-favorable operation, and evaluate offered-population
  SLO attainment, tail latency, background utility, fairness, route diversity,
  and control overhead on an actual four-node Perlmutter deployment.
  \item We provide a fresh held-out native result, paper-policy comparison,
  and source-bound hierarchy receipt with explicit claim boundaries and
  preserved failure artifacts.
\end{enumerate}

\section{Background and Motivation}

\subsection{Disaggregation moves the bottleneck}

A request with prompt length $n$ and output length $m$ can use one of two broad
paths.  A remote path computes prefill on $P_i$, transfers the KV cache, and
decodes on $D_j$.  Its TTFT includes prefill queueing and service, KV movement,
receiver queueing, and the first decode step.  A local-deflection path performs
chunked prefill on $D_j$ and removes inter-node KV movement, but each chunk
inflates a mixed prefill/decode step.  Thus

\begin{align}
\widehat{T}_{\mathrm{remote}} &= T_{P_i}^{q+s} + T_{i\rightarrow j}^{KV}
                               + T_{D_j}^{q+1}, \\
\widehat{T}_{\mathrm{local}}  &= T_{D_j}^{q} + T_{D_j}^{\mathrm{chunk+decode}}.
\end{align}

The hidden coupling is that $T_{i\rightarrow j}^{KV}$ depends on source and
receiver traffic, transport inflight work, NIC state, and competing
communication, while the local term depends on active sequences, KV occupancy,
and the TBT slack of existing decodes.  A decision made from one term can be
correct at admission and harmful by execution unless resources and lifecycle
are reserved consistently.

\subsection{Why shared HPC is a distinct setting}

Perlmutter GPU nodes contain four A100 GPUs and four Cassini Slingshot 11 NICs,
each rated at 200 Gb/s (25 GB/s), on a three-hop dragonfly
fabric~\cite{perlmutter,nerscnetwork}.  This gives substantial aggregate
bandwidth but also several independently changing contention domains: GPU
compute and HBM, PCIe and host service, NIC traffic classes, inter-group links,
and collective communication.  A scheduler cannot infer the active bottleneck
from ``local'' or ``remote'' alone.

The production objective is also not one scalar latency.  An interactive
tenant may require protected TTFT and end-to-end service; batch/background work
must retain bounded service; and a failure should release every reservation.
We therefore evaluate good requests out of the \emph{offered} population,
tenant fairness, and terminal outcomes, not only quantiles among completions.

\subsection{Motivating failure of partial policies}

Our held-out workload contains a decoder-hot regime and a P-only
remote-favorable regime.  The former overloads decode service with simultaneous
local and remote pressure.  The latter pre-seeds exact 4,094-token source hits
and loads both decoders, making remote source choice valuable.  No single
static path wins both.  In the final fresh run, the strongest fixed and
predictor policies exhibit 9.8--50.6 second stressed p99 latency, whereas a
joint policy stays near 3.3 seconds.  This gap is not produced by an isolated
network microbenchmark; it arises from the interaction of offered requests,
actual model execution, KV movement, decoder service, and admission.

\section{System Model and Objective}

Let $\mathcal{P}$ and $\mathcal{D}$ be prefill and decode instances.  A request
$r$ has tenant, arrival, deadline, prompt/cache identity, output demand, and a
candidate set $\mathcal{E}_r$.  A candidate
$e=(p,d,\rho)$ chooses source $p$, decoder $d$, and route $\rho\in
\{\mathrm{local},\mathrm{remote}\}$.  Each candidate carries a resource vector

\begin{equation}
  \begin{aligned}
  \mathbf{w}_e = (&w_{\mathrm{decode}}, w_{\mathrm{local\ prefill}}, w_{KV},\\
                   &w_{\mathrm{remote\ prefill}}, w_{\mathrm{semantic}}).
  \end{aligned}
\end{equation}

The controller consumes an atomic telemetry batch $\mathbf{s}$ containing
vLLM running/waiting requests and KV occupancy, cache-residency evidence,
LMCache semantic/byte inflight state, supported Cassini signals, endpoint
health, and controller-owned reservations.  Unsupported signals are unknown,
not zero.

Conceptually, \system chooses a feasible candidate that minimizes predicted
service plus externality and business debt:

\begin{equation}
 \begin{aligned}
 e^* = \arg\min_{e\in\mathcal{E}_r\cap\mathcal{F}(\mathbf{s})}
  \{\widehat{T}_e &+ \lambda_f C_e^{fabric}
  + \lambda_q C_e^{service}\\
  &+ \lambda_b C_{r,e}^{business}\}.
 \end{aligned}
\end{equation}

$\mathcal{F}$ enforces cache, capacity, freshness, deadline, health, and
transport guards.  The implementation uses calibrated service priors,
resource envelopes, tenant reservations, and causal receipts rather than
assuming that this expression is an exact latency oracle.  Requests can be
admitted, deferred in a bounded queue, or rejected.  Every admitted request
owns a lease that is released on first response, completion, explicit route
failure, timeout, or abort.

\section{Design}

\subsection{Identity-bound state plane}

Node agents export scheduler and endpoint state; pair agents bind it to model,
engine, topology, epoch, and monotonically increasing sequence identities.  A
request-triggered collector obtains one all-pair batch with bounded timeout and
one validation-only retry.  The coordinator atomically installs only a complete
and causally valid batch.  A stale fallback is permitted only within a
tenant-specific grace and still passes resource and health guards.

Cross-layer signals retain five fields: value, unit, support status, source,
and sample identity.  This matters on managed HPC systems where some counters
are available through user-visible sysfs or endpoint APIs and others require
operator integration.  Fabric-unavailable must not silently become
fabric-idle.

\subsection{Global candidate and atomic commit}

For each request, \system constructs feasible local and remote candidates over
the physical $P\times D$ mesh.  Cache identity determines whether a source can
supply a full hit, partial reuse, or a confirmed miss.  The coordinator scores
candidates, admits or queues the request, and returns an immutable commit that
contains route, source, decoder, edge, telemetry sequence, resource bindings,
and optional joint actuation.

The frontend reserves decoder placement from that commit and sends a remote
request to the committed prefill source.  The pair router independently checks
the commit before invoking vLLM/LMCache.  A mismatch fails closed.  This
separation prevents the common evidence bug of confusing the committed decode
destination with the prefill source on a cross-pair edge.

\subsection{Business admission and fairness}

Protected tenants receive bounded decoder admission and service-lane capacity;
background tenants receive explicit limits, maximum wait, and starvation
escape.  Fairness and capacity are part of candidate feasibility before
latency tie-breaking.  A request that cannot be safely admitted appears as a
global reject in the offered-population metric rather than disappearing from a
completed-only distribution.

The policy also tracks controller-owned virtual service.  Among remote
candidates whose calibrated scores differ only within telemetry uncertainty
and that have equivalent cache and business semantics, \system selects the
source and edge with lower virtual finish.  It never forces a materially worse
edge to satisfy a quota.  This breaks deterministic source ties while retaining
the original service optimum and yields causal source-balance receipts.

\subsection{Cross-layer actuation}

An admitted commit can carry route-specific resource limits, transfer semantic
concurrency, receiver staggering, and a shared-fabric budget.  Critical
transport or endpoint failures remain hard guards; non-critical pressure is
priced against the resource envelope.  Local deflection is available when
remote movement is harmful, while remote execution remains available when
decoder-local prefill would consume scarce service.

This mechanism is intentionally broader than a network penalty.  The same
decision can respond to decoder running-set inflation, a full remote cache hit,
fabric pressure, protected business demand, and pair health.  The action and
the signals that justified it are recorded together.

\subsection{Hierarchical fan-in}

At large logical pair counts, pair agents locally rank their route candidates.
Shard reducers forward a bounded frontier and an omission receipt; a global
coordinator retains final policy ownership.  The request identity, telemetry
epoch, omitted-pair count, and candidate fingerprint cross all levels.  This
bounds wire payload and global candidate count, but it does not by itself prove
native inference scalability; Section~\ref{sec:scale} keeps that boundary
explicit.

\section{Implementation}

The prototype comprises approximately four functional planes.

\tightparagraph{Serving carrier}
Four vLLM 0.26.0 engines run Qwen2.5-7B-Instruct at tensor parallelism four.
Two act as prefill producers and two as decode consumers.  The common remote
path is LMCacheConnectorV1 with NIXL/UCX over the CXI provider.  Local execution
uses vLLM chunked prefill on the decoder.  The experiment disables decoder
prefix caching so LMCache source evidence and local work are not conflated.

\tightparagraph{State and control}
The Python implementation contains strict dataclasses for resource vectors,
pair telemetry, cross-layer signals, route candidates, decisions, failure
reports, and hierarchy receipts.  An asynchronous coordinator joins telemetry
collection, global admission, stream feedback, and release.  Frontend and pair
router APIs preserve request IDs and commitment hashes across HTTP and SSE
boundaries.

\tightparagraph{Perlmutter instrumentation}
Cassini readers expose per-NIC and traffic-class counters with explicit support
classification.  vLLM scheduler snapshots supply running, waiting, and KV
occupancy.  LMCache supplies semantic and byte-inflight state.  NCCL observer
schemas exist and are exercised in opt-in co-job tests, but NCCL latency signals
were unavailable in the final held-out headline and are reported as such.

\tightparagraph{Failure discipline}
The system rejects wrong epoch, topology, sequence, hash, stale state, missing
candidate ownership, duplicate lifecycle events, and leaked leases.  The
artifact retains unsuccessful launches and validator failures.  Experiments run
inside a four-hour no-shell interactive allocation; compute occurs only in
foreground Slurm steps.

\section{Evaluation Methodology}

We ask five questions:

\begin{enumerate}
  \item Does \system improve offered-population SLO attainment and tail latency
  under decoder-hot contention?
  \item Does it retain remote execution and use the physical mesh when exact
  source reuse makes remote favorable?
  \item Does foreground protection preserve bounded background utility and
  tenant fairness?
  \item How does it compare with paper-derived network-aware selection and
  prefill deflection on the same actual carrier?
  \item Does the hierarchy bound global fan-in as logical pair count grows?
\end{enumerate}

\subsection{Testbed and workload}

The native testbed is one Perlmutter allocation with four GPU nodes and 16 A100
GPUs.  Each node contributes one TP4 engine.  Perlmutter uses Slingshot 11 with
four Cassini NICs per GPU node~\cite{perlmutter}.  We use an 8,192-token model
limit, 16 maximum sequences per engine, 3,000 ms TTFT, 150 ms TBT/TPOT, and
8,000 ms end-to-end SLOs.

The held-out contract fixes seed 20260825, deterministic sub-spacing arrival
jitter, block order, prompt namespace transform, request counts, profile and
source hashes before the fresh execution.  The three analysis regimes are:

\begin{itemize}
  \item \textbf{normal}: 60 foreground requests across two control blocks;
  \item \textbf{miss-hot}: 120 foreground requests plus simultaneous remote
  and decoder-local background pressure across both decode destinations;
  \item \textbf{remote-favorable}: 30 foreground requests with exact
  4,094-token P-only source hits while 1,344 background requests load both
  decoders.
\end{itemize}

We report offered, completed, and SLO-good requests; TTFT, TPOT and end-to-end
p50/p95/p99 among completions; route and edge counts; background terminal
outcomes; per-tenant completion and Jain fairness; telemetry overhead and
support; and all failure/rejection counts.

\subsection{Baselines}

The independently validated matrix contains local and remote fixed policies,
a calibrated predictor, and queue/GPU observe-only routing.  All use the same
model, engines, request population, and remote connector.

The C10 extension adds two paper-policy reproductions.  \textbf{NetKV} uses
remote candidates and minimizes KV transfer plus decoder queue and first-step
cost, using 25 GB/s per TP shard, supported Cassini congestion, LMCache
self-inflight, and vLLM scheduler state~\cite{netkv}.  \textbf{Kairos-$X512$}
uses $\alpha=1.3$, a 0.9 TBT safety margin, and an actual 512-token decoder
batching ceiling~\cite{kairos}.  Stock vLLM lacks Kairos's per-request dynamic
chunk schedule, and the reported repository was unavailable at our freeze
time; this is therefore a declared $X=\{512\}$ subset, not the authors' full
implementation.

Both paper baselines disable \system's business priority, fairness reservation,
pair scaling, shared-fabric budget, and stagger.  A compatibility adapter emits
frozen-workload lifecycle evidence but never waits or throttles.  The analyzer
verifies \texttt{policy\_effect=none} and \texttt{wait\_ns=0} for every receipt.

\subsection{Claim boundaries}

\begin{center}
\fbox{\parbox{0.94\columnwidth}{
The C9 headline is a fresh-allocation, one-shot, source-bound independent
validation.  C10 is an actual-system, same-population \emph{post-hoc}
extension because its adapters were frozen after that allocation began.  The
hierarchy experiment is CPU control-plane evidence only.  We do not claim a
full Kairos reproduction, an independent C10 SOTA result, native scale beyond
four nodes, or supported NCCL latency telemetry in the held-out headline.
}}
\end{center}

\section{Evaluation}

\subsection{Fresh held-out native result}

Table~\ref{tab:headline} reports \system on all three regimes.  Every offered
foreground request completes within all SLOs.  The remote-favorable phase uses
29 remote requests and one local request while retaining a 3.36-second p99.

\begin{table}[t]
\centering
\caption{Fresh four-node C9 held-out result.  Counts are
offered/completed/SLO-good.}
\label{tab:headline}
\small
\begin{tabular}{lrrr}
\toprule
Regime & Count & E2E p50 & E2E p99 \\
\midrule
Normal & 60/60/60 & 2.910 s & 3.150 s \\
Miss-hot & 120/120/120 & 3.058 s & 3.337 s \\
Remote-favorable & 30/30/30 & 3.067 s & 3.357 s \\
\bottomrule
\end{tabular}
\end{table}

Table~\ref{tab:reductions} shows relative p99 reduction.  The strongest fixed
arm is selected per stressed regime by the frozen analyzer, rather than chosen
after seeing \system.  The improvement is large because partial policies
either queue behind the hot decoder, retain an expensive remote path, or fail
to coordinate admission with execution.

\begin{table}[t]
\centering
\caption{\system E2E p99 reduction relative to matched baselines.}
\label{tab:reductions}
\small
\begin{tabular}{lrr}
\toprule
Baseline & Miss-hot & Remote-favorable \\
\midrule
Strongest fixed & 66.12\% & 90.55\% \\
Predictor & 61.69\% & 93.36\% \\
Queue/GPU-only & 58.48\% & 93.54\% \\
\bottomrule
\end{tabular}
\end{table}

\subsection{Background utility and control cost}

Foreground protection does not delete the offered background population.  In
miss-hot blocks, background completes 1,204 of 1,404 requests (85.75\%); the
minimum block-by-tenant completion is 76.50\%, the two-tenant Jain fairness is
0.9979, and service-lane failures are 0.93\%.  In remote-favorable, background
completes 1,344/1,344.  These outcomes include explicit rejects and failures in
the denominator.

Request-triggered telemetry collection has 28.62 ms p50 and 132.42 ms p99;
admission has 29.46 ms p50 and 133.26 ms p99.  All 30 remote-favorable decisions
carry explicit support classifications.  Cassini signals are supported for 29
and LMCache inflight signals for all 30.  The source/edge virtual-service
binding appears on 29 decisions.  Unavailable NCCL and LMCache latency signals
remain unavailable.

\subsection{Physical mesh use}

Across regimes, \system uses both local decoder edges and all four remote
edges: $P_0\!\rightarrow\!D_0$, $P_0\!\rightarrow\!D_1$,
$P_1\!\rightarrow\!D_0$, and $P_1\!\rightarrow\!D_1$.  This rules out a
degenerate ``global'' controller that always sends requests to one source or
one receiver.  Every cross-pair remote completion carries a commit whose
source, destination, edge identity, cache evidence, and telemetry sequence are
consistent.

\subsection{Paper-policy comparison}

Table~\ref{tab:sota} compares actual-system C10 results.  NetKV is strong in the
normal phase, where it completes all requests and trails \system p99 by 4.54\%.
Under miss-hot contention it completes 117/120 but only 73 meet the SLO; under
remote-favorable load it admits one request, whose 53.9-second latency misses
the SLO.  \system reduces NetKV p99 by 75.79\% and 93.77\% in the two stressed
regimes and increases SLO-good by 47 and 30.

Kairos-$X512$ completes half of normal requests, rejects every miss-hot victim,
and completes one remote-favorable victim.  Because miss-hot has zero
completions, its latency quantiles are undefined.  We report a 120-request
completion/SLO advantage, not a fabricated infinite latency speedup.  On the
defined remote tail, \system reduces p99 by 30.82\%.

\begin{table*}[t]
\centering
\caption{Actual four-node C10 paper-policy comparison.  SLO is good/offered;
latency is completion E2E p99.  Kairos is the declared $X=\{512\}$ subset.}
\label{tab:sota}
\small
\begin{tabular}{lrrrrrr}
\toprule
& \multicolumn{2}{c}{Normal} & \multicolumn{2}{c}{Miss-hot}
& \multicolumn{2}{c}{Remote-favorable} \\
Policy & SLO & p99 & SLO & p99 & SLO & p99 \\
\midrule
\system & 60/60 & 3.150 s & 120/120 & 3.337 s & 30/30 & 3.357 s \\
Kairos-$X512$ & 30/60 & 3.223 s & 0/120 & undefined & 1/30 & 4.853 s \\
NetKV reproduction & 60/60 & 3.300 s & 73/120 & 13.780 s & 0/30 & 53.914 s \\
\bottomrule
\end{tabular}
\end{table*}

The comparison demonstrates the motivating point, not that the referenced
systems lack value.  Network-aware remote selection cannot rescue a request
when the correct action is decoder-local prefill or business-aware admission;
a fixed safe deflection subset cannot exploit all time-varying decode slack or
cache/fabric alternatives.  \system retains both action families and commits
them with shared state.

\subsection{Control-plane scale}
\label{sec:scale}

We replay the same two-route candidate population at 2, 8, 32, 128, 512, and
1,024 logical pairs for 15 repetitions.  Pair ownership is prepartitioned, so
pair-agent timing measures local ranking and receipt construction rather than
repeatedly scanning the global population.

At 1,024 pairs, the bounded path forwards 256 of 2,048 candidates and records
896 omitted pairs.  Serialized global payload falls from 666,815 to 83,358
bytes (87.50\%).  Full scan p50/p99 is 49.74/152.84 ms; pair-agent p50 is
29.65 ms; bounded-global p50 is 55.43 ms; and bounded end-to-end p50/p99 is
85.33/158.24 ms.  Thus global payload and candidate fan-in are bounded, while
this single-process Python replay is not faster than full scan.  Distributed
wall-clock and native utility beyond four nodes remain future work.

\section{Failure Analysis and Reproducibility}

The final result follows hundreds of preserved development versions rather
than hiding unsuccessful candidates.  Several route-only and threshold
policies improved completed-only latency by rejecting too much work; they were
discarded by offered-population gates.  Other runs exposed stale telemetry,
source/destination evidence aliasing, and frontend lifecycle failures.

C10 itself produced two non-performance failures.  The first baseline adapter
removed a required evidence receipt while correctly removing policy.  The fix
emits a no-effect lifecycle receipt and the analyzer proves zero wait and no
throttle.  The second failure occurred after NetKV executed a valid
$P_0\!\rightarrow\!D_1$ request: a legacy validator treated the committed
decoder destination as the prefill source.  A mesh-aware adapter now validates
commit source, decoder, and edge separately without changing the workload or
score.  The successful Kairos result was hash-bound and reused; only NetKV was
executed under the corrected evidence source.

The artifact includes run contracts, source inventories, raw request and
router ledgers, endpoint evidence, logs, failed result roots, analyses, and
SHA-256 digests.  Independent and post-hoc results are stored separately.

\section{Related Work}

\tightparagraph{LLM serving and disaggregation}
vLLM reduces KV-memory waste and enables flexible batching with
PagedAttention~\cite{vllm}.  DistServe separates prefill and decode to optimize
goodput under phase-specific SLOs~\cite{distserve}.  Mooncake builds a
KV-centric disaggregated cache and global scheduler at production
scale~\cite{mooncake}; LMCache exposes reusable, movable KV state across
engines~\cite{lmcache}.  \system uses vLLM and LMCache as its actual carrier and
focuses on a shared-HPC control loop that joins their state with fabric,
topology, failure, and tenant policy.

\tightparagraph{Network-aware routing}
NetKV introduces a network cost oracle and an $O(|D|)$ decode-instance selector
that combines KV transfer and decoder cost~\cite{netkv}.  Its published
evaluation uses a 64-GPU four-tier fat-tree simulator driven by Mooncake traces.
\system extends the action space to local prefill, admission, business
fairness, and joint resource commitment, and evaluates a NetKV policy
reproduction on an actual Slingshot carrier.  MRC provides production-grade
multipath, congestion control, loss recovery, and path resilience at the
transport layer~\cite{mrc}.  Such transport mechanisms can improve the
available paths beneath \system but do not replace inference-semantic and
business orchestration.

\tightparagraph{Prefill deflection}
Kairos estimates regular prefill TTFT and searches decode-node chunk schedules
that preserve TBT SLO, reporting large P95 TTFT and SLO gains on a 2P2D A100
cluster~\cite{kairos}.  \system likewise includes decoder-local chunked
prefill, but decides it jointly with remote cache/fabric candidates and global
admission.  Our $X512$ experiment is deliberately narrower than Kairos and is
used only as a same-carrier mechanism comparison.

\section{Limitations}

The independently validated native experiment uses four nodes and one model.
It proves a four-node Perlmutter mechanism and performance result, not
facility-wide production superiority.  The 1,024-pair result is CPU-only and
does not measure distributed agent wall-clock, NIC traffic, GPU inference, or
failure convergence.

The final held-out run has supported Cassini and LMCache inflight state but not
NCCL collective latency or LMCache transfer-tail signals.  Those schemas and
co-job mechanisms exist, yet a future source-frozen experiment must make them
supported and perform a causal ablation before claiming that those signals
drive the headline gain.

C10 is post-hoc and must be repeated unchanged on a fresh allocation for an
independent SOTA-extension claim.  Kairos-$X512$ is not full Kairos.  Results
may also change for larger models, longer contexts, more cache tiers, different
topologies, and operator-level scheduling privileges.  At user scope, \system
controls only its allocation and treats other jobs as exogenous; it does not
modify facility policy or other users' workloads.

\section{Conclusion}

Disaggregated LLM inference on a shared HPC fabric is not a binary local-versus-
remote routing problem.  Decoder service, KV movement, receiver load, fabric
state, cache identity, tenant utility, and failures form one coupled lifecycle.
\system makes that lifecycle the unit of orchestration.  On an actual
four-node Perlmutter deployment, a fresh held-out run sustains 100\% foreground
SLO attainment and roughly 3.3-second stressed p99 while matched partial
policies degrade to 8--53 seconds or reject work.  Background and telemetry
gates show that the gain is not obtained by silently dropping the rest of the
population.  Post-hoc paper-policy and logical-scale results strengthen the
motivation while preserving their narrower claim boundaries.  The next step is
an unchanged independent paper-policy rerun and native scale beyond four nodes,
not another threshold search on the proven workload.

\bibliographystyle{plain}
\bibliography{references}

\end{document}
